\chapter{Nonconjugate block integration and numerical certification}
\label{ch:nonconj}
\label{app:pg-mf-derivation}
\label{app:variational-bounds}
\chapterpromise{Laplace, variational, expectation-propagation, Pólya--Gamma, and quadrature calculations for segment marginal likelihoods, with conditional propagation of segmentwise log-score error through the partition recursions.}

\begin{statusbox}[Exact global algebra, approximate local target]
Once an approximate block array is supplied, the partition recursions are evaluated exactly for
that array. This is not the same as exact inference under the reference nonconjugate model. The gap
is controlled only when a local certificate verifies Assumption~\ref{ass:uniform-block-error} on
all blocks reachable by the declared partition support.
\end{statusbox}
\label{sec:nonconj}

\paragraph{Plate model (non-conjugate single segment with PG augmentation).}
The generative graph for a non-conjugate block $(i,j]$ is identical to Figure~\ref{fig:plate-ef}
with $\theta\sim\pi(\theta)$ replaced by a non-conjugate prior (e.g., Gaussian on a log or logit
scale). For Pólya--Gamma-augmented binomial logistic blocks (Theorem~\ref{thm:pg}), introducing
the latent scale-mixture variables $\omega_t\sim\mathrm{PG}(m_t,0)$ restores conditional
conjugacy and is reflected in the plate as a second observed-indexed latent:
\begin{equation}
\theta\sim\mathcal{N}(\nu,\rho^2),\quad
\omega_t\stackrel{\mathrm{ind}}{\sim}\mathrm{PG}(m_t,0),\quad
y_t\mid\theta,\omega_t,m_t\stackrel{\mathrm{ind}}{\sim}\mathrm{Binom}(m_t,\sigma(\theta))\text{ with \eqref{eq:pg-identity}},
\label{eq:plate-nonconj}
\end{equation}
where $\sigma$ is the logistic function.
\begin{figure}[H]
\centering
\begin{tikzpicture}
  \node[obs] (y) {$y_t$};
  \node[latent, above=0.9cm of y] (om) {$\omega_t$};
  \node[latent, left=1.2cm of om] (th) {$\theta$};
  \node[const, left=0.9cm of th] (nr) {$\nu,\rho^2$};
  \node[const, right=0.9cm of y] (m) {$m_t$};
  \edge {nr} {th}
  \edge {th,om,m} {y}
  \edge {m} {om}
  \plate {p1} {(y)(om)} {$t\in(i,j]$}
\end{tikzpicture}
\caption{Plate model for a non-conjugate block with Pólya--Gamma augmentation. Conditional on $\omega_t$, the block is quadratic in $\theta$ with a Gaussian posterior (Proposition~\ref{prop:pg-mf}); the same template applies to Laplace and variational routines with the $\omega_t$ node omitted.}
\label{fig:plate-nonconj}
\end{figure}

In many applications the single-segment parameterization is not conjugate to the observation
model---for example, log-normal priors for positive means, heavy-tailed shrinkage priors, or
logistic/negative-binomial links with structural offsets. This section develops a plug-in
toolkit for computing approximate block evidences $A^{(0)}_{ij}$ and moments $A^{(r)}_{ij}$
when conjugacy is unavailable, while keeping the dynamic-programming layer of
Section~\ref{sec:dp} unchanged. When approximate block evidences are used, the DP is exact for
the approximate block-score model; additional error statements require explicit blockwise
approximation control.

We focus on one-dimensional (or low-dimensional) canonical GLMs, where for a block $(i,j]$ we
must approximate integrals of the form
\[
A^{(0)}_{ij}
 = \int \Big[\prod_{t=i+1}^j p(y_t\mid\theta)\Big] p(\theta)\,d\theta,\qquad
A^{(r)}_{ij}
 = \int \Big[\prod_{t=i+1}^j p(y_t\mid\theta)\Big] p(\theta)\,m_\star(\theta)^r\,d\theta,
\]
where $m_\star(\theta)=m(\theta;d_\star)$ is the declared reporting functional at a fixed reference descriptor. We describe four complementary approaches:
(i) first-order Laplace approximation; (ii) variational lower bounds (Jaakkola--Jordan \citep{jaakkola2000logisticvb} type
for logistic regression and mean-field for more general GLMs); (iii) expectation propagation
(EP); and (iv) P\'olya--Gamma augmentation for logistic and negative-binomial blocks. Throughout
we use the family-specific descriptor $d_t$ and any separately declared power weight $a_t$ and cumulative statistics $(S_{ij},W_{ij},H_{ij})$.

\subsection{Model and notation for GLM blocks}

Consider a one-parameter canonical GLM on a block $(i,j]$:
\[
\log p(y_t\mid\theta)
 = w_t\big\{ y_t\theta - b(\theta)\big\} + \log h(y_t,w_t),
\]
where $b$ is the cumulant function and $h$ collects base measures (including any $w_t$-dependent
terms). The block log-likelihood is
\[
\ell_{ij}(\theta)
 := \sum_{t=i+1}^j w_t\{y_t\theta - b(\theta)\} + H_{ij},
\qquad
H_{ij}:=\sum_{t=i+1}^j \log h(y_t,w_t).
\]
Let $\pi(\theta)=p(\theta)$ be an arbitrary prior (not necessarily conjugate), and write
$\Psi_{ij}(\theta):=\ell_{ij}(\theta)+\log\pi(\theta)$ for the block log-posterior.

Derivatives are
\[
\ell_{ij}'(\theta) = S_{ij} - W_{ij} b'(\theta),\qquad
\ell_{ij}''(\theta) = -W_{ij} b''(\theta),
\]
and
\[
\Psi_{ij}'(\theta) = \ell_{ij}'(\theta) + (\log\pi)'(\theta),\qquad
\Psi_{ij}''(\theta) = \ell_{ij}''(\theta) + (\log\pi)''(\theta).
\]

\subsection{Existence and uniqueness of block modes}

\begin{lemma}[Strict log-concavity and uniqueness]
\label{lem:concave-unique}
Suppose $b$ and $\log\pi$ are twice differentiable on a convex parameter domain,
$b''(\theta)>0$ for all $\theta$ (canonical GLM), the prior is log-concave
$((\log\pi)''(\theta)\le 0)$, and $W_{ij}>0$. Then $\Psi_{ij}$ is strictly concave. If, in
addition, $\Psi_{ij}$ is proper and attains its supremum on the parameter domain (for example by
coercive tails or compact-domain truncation), the maximizer $\hat\theta_{ij}$ exists and is unique.
\end{lemma}

\begin{proof}
Since $W_{ij}>0$ and $b''(\theta)>0$, we have $-W_{ij}b''(\theta)<0$. Adding
$(\log\pi)''(\theta)\le 0$ yields $\Psi_{ij}''(\theta)<0$ wherever the derivatives exist, so
$\Psi_{ij}$ is strictly concave. Strict concavity gives uniqueness of any maximizer; existence is
provided by the stated attainment assumption.
\end{proof}

\subsection{Laplace approximation for single-block evidence}


\begin{theorem}[First-order Laplace approximation with a uniform remainder]
\label{thm:laplace}
Index a sequence of blocks by an effective information size $N=W_{ij}\to\infty$. Let
$\widehat\theta_N$ be the unique interior maximizer of $\Psi_N$ and
$H_N=-\Psi_N''(\widehat\theta_N)$. Assume there are constants $c,C,\delta>0$, independent of the
block, such that: (i) $cN\le-\Psi_N''(\theta)\le CN$ throughout
$|\theta-\widehat\theta_N|\le\delta$; (ii) $\Psi_N$ is four times continuously differentiable on
that neighborhood and $|\Psi_N^{(r)}(\theta)|\le CN$ for $r=3,4$; and (iii) the integral outside
the neighborhood is at most the local Gaussian leading term times $O(N^{-1})$. Then
\begin{equation}
 A^{(0)}_{ij}=e^{\Psi_N(\widehat\theta_N)}
 \left(\frac{2\pi}{H_N}\right)^{1/2}\{1+O(N^{-1})\},
\label{eq:laplace-A0-mult}
\end{equation}
and hence
\begin{equation}
 \log A^{(0)}_{ij}=\Psi_N(\widehat\theta_N)+\tfrac12\log(2\pi)-\tfrac12\log H_N+O(N^{-1}).
\label{eq:laplace-A0}
\end{equation}
For a bounded twice continuously differentiable test function $g$ whose first two derivatives are uniformly bounded in the same neighborhood,
\begin{equation}
 \frac{M^{[g]}_{ij}}{A^{(0)}_{ij}}=g(\widehat\theta_N)+O(N^{-1}),
\qquad M^{[g]}_{ij}=\int g(\theta)e^{\Psi_N(\theta)}\,d\theta.
\label{eq:laplace-moment}
\end{equation}
\end{theorem}
\begin{proof}
Set $u=\sqrt{H_N}(\theta-\widehat\theta_N)$. Taylor's theorem through fourth order gives, on the
local neighborhood,
\[
 \Psi_N(\theta)-\Psi_N(\widehat\theta_N)
 =-\tfrac12u^2+a_Nu^3N^{-1/2}+b_N(u)u^4N^{-1},
\]
where $a_N$ and $b_N(u)$ are uniformly bounded by assumptions (i)--(ii). On a growing central
region $|u|\le N^{1/12}$, expand the exponential. The term linear in $u^3$ integrates to zero
against the symmetric Gaussian kernel; the remaining terms are bounded by an integrable Gaussian
polynomial times $N^{-1}$. The part of the local neighborhood outside this central region is
exponentially smaller by the quadratic term, and assumption (iii) controls the nonlocal tail.
After the Jacobian $H_N^{-1/2}$, this proves the relative $1+O(N^{-1})$ expansion. Applying the
same argument to the numerator after the second-order Taylor expansion
$g(\widehat\theta_N+u/\sqrt{H_N})=g(\widehat\theta_N)+g'(\widehat\theta_N)u/\sqrt{H_N}+O(u^2/N)$
and again using Gaussian symmetry yields the ratio statement.
\end{proof}


\paragraph{Efficient Newton updates using block summaries.}
For canonical GLMs,
\[
\Psi_{ij}'(\theta) = S_{ij} - W_{ij} b'(\theta) + (\log\pi)'(\theta),\quad
\Psi_{ij}''(\theta) = -W_{ij} b''(\theta) + (\log\pi)''(\theta),
\]
so Newton iteration is
\[
\theta^{(m+1)} \leftarrow \theta^{(m)}
 - \frac{S_{ij} - W_{ij} b'(\theta^{(m)}) + (\log\pi)'(\theta^{(m)})}
        {-\,W_{ij} b''(\theta^{(m)}) + (\log\pi)''(\theta^{(m)})},
\]
with cost independent of block length once $S_{ij},W_{ij}$ are available.

\paragraph{Implementation.}
Algorithm~\ref{alg:laplace} implements the Laplace block-evidence approximation using a Newton
or quasi-Newton solver for the block mode and the corresponding Hessian correction. The
approximation is local but often highly accurate in moderate-to-large blocks. The remaining
non-conjugate routines (Algorithms~\ref{alg:jj}--\ref{alg:precompute-nonconj}) follow the same
template: a block-local optimizer feeds an approximate log-evidence into the DP.

\begin{algorithm}[H]
\caption{GLMBlock\_Laplace$(y,w,i,j,\texttt{b},\texttt{prior})$}
\label{alg:laplace}
\KwIn{$y_{i+1:j}$, weights $w_{i+1:j}$; cumulant $b(\theta)$; prior $\pi(\theta)$}
Compute $S\leftarrow \sum_{t=i+1}^j w_t y_t$;\quad $W\leftarrow \sum_{t=i+1}^j w_t$;\quad $H\leftarrow \sum_{t=i+1}^j \log h(y_t,w_t)$\;
Define $\Psi(\theta)=H + S\theta - W b(\theta) + \log \pi(\theta)$ and its derivatives\;
Initialize $\theta^{(0)}$ (e.g., solve $b'(\theta)=S/W$ or use prior mean)\;
\For{$m=0,1,\dots$ until convergence}{
  $g\leftarrow S - W b'(\theta^{(m)}) + (\log\pi)'(\theta^{(m)})$\;
  $h\leftarrow - W b''(\theta^{(m)}) + (\log\pi)''(\theta^{(m)})$\;
  \uIf{$|g|/|h|<\varepsilon_{\text{newton}}$}{break}
  $\theta^{(m+1)} \leftarrow \theta^{(m)} - g/h$\;
}
$\hat\theta\leftarrow \theta^{(m)}$;\quad $H^\star\leftarrow -\,\Psi''(\hat\theta)$\;
$\log \widehat{A}^{(0)}_{ij}\leftarrow \Psi(\hat\theta) + \tfrac{1}{2}\log(2\pi) - \tfrac{1}{2}\log H^\star$\;
\KwOut{$\widehat{A}^{(0)}_{ij}$ and (optional) moments via \eqref{eq:laplace-moment}}
\end{algorithm}

\subsection{Variational lower bounds for GLM blocks}

When a deterministic lower bound is desirable, introduce a variational distribution $q(\theta)$
and apply Jensen:
\[
\log A^{(0)}_{ij}
=\log \int q(\theta)\,\frac{e^{\Psi_{ij}(\theta)}}{q(\theta)}\,d\theta
\ \ge\
\int q(\theta)\,\Psi_{ij}(\theta)\,d\theta + \mathcal{H}[q]
 \;:=\; \mathcal{L}_{ij}(q).
\]

\emph{Jaakkola--Jordan bound for logistic blocks.}
For Bernoulli/logistic likelihoods, a standard quadratic bound yields a tractable Gaussian
variational posterior. One convenient form (for sigmoid $\sigma$) is:
\[
\log\sigma(z)\ \ge\ \log\sigma(\xi)+\frac{z-\xi}{2}-\lambda(\xi)(z^2-\xi^2),
\qquad
\lambda(\xi)=\frac{\tanh(\xi/2)}{4\xi},
\]
with variational parameter $\xi>0$. Applying this bound to the appropriate signed linear
predictors yields a quadratic surrogate in $\theta$, so the optimal Gaussian $q(\theta)$ and
$\xi$ admit closed-form coordinate updates.

\begin{proposition}[Monotone improvement of variational block bounds]
\label{prop:var-monotone}
Every exact coordinate-ascent update of the local variational parameters (for example $\xi$ in a
logistic block) or of a Gaussian $q(\theta)$ leaves the block evidence lower bound
$\mathcal L_{ij}(q)$ non-decreasing. If the iterates converge to an interior fixed point and the
objective is differentiable there, that fixed point is coordinatewise stationary. At every
iteration, $\underline{A}^{(0)}_{ij}:=\exp\{\mathcal{L}_{ij}(q)\}$ is a certified lower bound on
$A^{(0)}_{ij}$.
\end{proposition}

\begin{proof}
Each coordinate update is an exact maximization of $\mathcal{L}_{ij}(q)$ with respect to a subset
of variables holding the others fixed, so the ELBO is non-decreasing. At a converged interior fixed
point, the first-order derivative along each updated coordinate block vanishes under the stated
regularity, which is coordinatewise stationarity. Because $\mathcal{L}_{ij}$ is a Jensen lower
bound on $\log A^{(0)}_{ij}$ at every iterate, exponentiating yields a valid evidence lower bound.
\end{proof}

\begin{algorithm}[H]
\caption{LogisticBlock\_JJ$(y,m,w,i,j,\texttt{prior})$ \quad (variational lower bound)}
\label{alg:jj}
\KwIn{Binomial counts $y_t\sim\mathrm{Binom}(m_t,p)$, logit link; prior $\pi(\theta)$}
Initialize $\xi$ and a Gaussian $q(\theta)=\mathcal{N}(\mu_q,\sigma_q^2)$\;
\Repeat{converged}{
  Update the quadratic bound parameter $\xi\leftarrow \sqrt{\mathbb{E}_q[\theta^2]}$\;
  Update Gaussian $q(\theta)$ by maximizing the block ELBO $\mathcal{L}_{ij}(q)$ (closed-form for quadratic surrogate)\;
}
Return $\underline{A}^{(0)}_{ij}=\exp\{\mathcal{L}_{ij}(q)\}$ and moments from $q$.
\end{algorithm}

\subsection{Expectation propagation (EP) for GLM blocks}
Expectation propagation (EP) \citep{minka2001ep} provides a deterministic approximation to the
block posterior and block evidence by iteratively refining local site approximations and matching
moments. In the formulation used here the prior is Gaussian, and each likelihood term
$p(y_t\mid\theta)$ is approximated by a Gaussian site
$\tilde t_t(\theta)=c_t\exp\{-\tfrac{1}{2}a_t\theta^2+b_t\theta\}$. After convergence, the product
of the Gaussian prior and sites is Gaussian and the approximate evidence is
$\widehat A^{(0,\mathrm{EP})}_{ij}=\int \pi(\theta)\prod_{t=i+1}^j \tilde t_t(\theta)\,d\theta$.
The constants $c_t$ are part of the returned log evidence; dropping them changes the DP score.
A non-Gaussian prior requires an additional site approximation or numerical integral and is not
covered by this Gaussian-posterior description. EP is neither a bound nor guaranteed to converge.

\begin{algorithm}[H]
\caption{GLMBlock\_EP$(y,w,i,j,\texttt{prior})$}
\label{alg:ep}
\KwIn{Canonical GLM; Gaussian prior $\pi(\theta)$; tolerance and maximum iterations}
Initialize site parameters $\{a_t,b_t,c_t\}_{t=i+1}^j$; set $\tilde p(\theta)\propto\pi(\theta)\prod_t \tilde t_t(\theta)$\;
\Repeat{moments stable or maximum iterations reached}{
  \For{$t=i+1$ \KwTo $j$}{
    Form cavity $\tilde p^{\setminus t}(\theta)\propto \tilde p(\theta)/\tilde t_t(\theta)$ and verify positive cavity precision\;
    Form tilted $p^{\text{tilt}}(\theta)\propto \tilde p^{\setminus t}(\theta)\,p(y_t\mid\theta)$\;
    Moment-match $p^{\text{tilt}}$ and update the site parameters $(a_t,b_t,c_t)$, damping if needed\;
  }
}
Return a convergence/failure flag, the approximate log evidence
$\log\widehat A^{(0,\mathrm{EP})}_{ij}=\log\int \pi(\theta)\prod_{t=i+1}^j c_t\exp\{-a_t\theta^2/2+b_t\theta\}\,d\theta$, and moments from $\tilde p$.
\end{algorithm}

An implementation should export the convergence flag and the maximum site residual together with the
block score. Downstream DP routines may then exclude failed EP blocks, retry them with stronger
damping, or mark the whole EP approximation as uncertified for the stability analysis.

\subsection{Pólya--Gamma augmentation for binomial-logistic blocks}
\begin{theorem}[P\'olya--Gamma identity and conditional conjugacy]
\label{thm:pg}
For binomial logistic blocks with counts $m_t$ and successes $y_t$, let $w_t\ge 0$ denote either
one or a declared likelihood-power/replication weight, define
$b_t=w_t m_t$ and $\kappa_t=w_t(y_t-\tfrac{m_t}{2})$, and introduce latent
$\omega_t\sim\mathrm{PG}(b_t,0)$. Then the weighted logistic likelihood term admits the Gaussian
scale-mixture representation
\begin{equation}
\left\{\frac{(e^\theta)^{y_t}}{(1+e^\theta)^{m_t}}\right\}^{w_t}
 = 2^{-b_t} e^{\kappa_t\theta}
   \int_0^\infty e^{-\omega_t\theta^2/2}\,p(\omega_t)\,d\omega_t.
\label{eq:pg-identity}
\end{equation}
The binomial coefficient, when present, is independent of $\theta$ and must be retained in the base-measure contribution to the block evidence. Conditionally on $\omega_{i+1:j}$, the remaining block likelihood is quadratic in $\theta$:
\begin{equation}
\log p(y_{i+1:j}\mid \theta,\omega)
 =
 -\tfrac{1}{2}\theta^2 \sum_{t=i+1}^j \omega_t
 + \theta \sum_{t=i+1}^j \kappa_t
 + \mathrm{const}.
\label{eq:pg-cond-quad}
\end{equation}
With a Gaussian prior $\pi(\theta)=\mathcal{N}(\nu,\rho^2)$, the conditional posterior
$p(\theta\mid y,\omega)$ is Gaussian.
\end{theorem}

\begin{proof}
\ProofStep{Step 1: the identity.}
Equation \eqref{eq:pg-identity} is the defining P\'olya--Gamma augmentation of
\citet{polson2013pg}: it expresses the binomial-logistic term as a Gaussian scale mixture in
$\theta$.

\ProofStep{Step 2: condition on $\omega$ and collect quadratic terms.}
Given $\omega_t$, each weighted term contributes $\kappa_t\theta-\tfrac12\omega_t\theta^2$ up to
constants, where the weight has already been absorbed into $b_t$ and $\kappa_t$. Summing over $t$
yields \eqref{eq:pg-cond-quad}.

\ProofStep{Step 3: combine with a Gaussian prior.}
Multiplying the quadratic-in-$\theta$ conditional likelihood by a Gaussian prior produces a
Gaussian conditional posterior, because Gaussian priors are conjugate to Gaussian likelihoods.
\end{proof}

\begin{proposition}[Mean-field variational approximation under P\'olya--Gamma]
\label{prop:pg-mf}
Under the same setup as Theorem~\ref{thm:pg} and a Gaussian prior $\pi(\theta)=\mathcal{N}(\nu,\rho^2)$,
the mean-field factorization $q(\theta,\omega)=q(\theta)\prod_{t=i+1}^j q(\omega_t)$ admits the
following closed-form coordinate-ascent updates:
(i) $q(\theta)=\mathcal{N}(\mu_q,\sigma_q^2)$ with
$\sigma_q^{-2}=\rho^{-2}+\sum_t \mathbb{E}_q[\omega_t]$ and
$\mu_q=\sigma_q^2(\rho^{-2}\nu+\sum_t \kappa_t)$;
(ii) $q(\omega_t)=\mathrm{PG}(b_t, c_t)$ with $c_t^2=\mathbb{E}_q[\theta^2]=\mu_q^2+\sigma_q^2$,
giving $\mathbb{E}_q[\omega_t]=\tfrac{b_t}{2 c_t}\tanh(c_t/2)$. These updates are monotone in the
block ELBO.
\end{proposition}

\begin{proof}
For a mean-field factor, the coordinate optimum satisfies
\[
\log q^*(\theta)
=\mathbb E_{q(\omega)}\!\left[\log p(y,\theta,\omega)\right]+\mathrm{const}.
\]
Combining Equation~\eqref{eq:pg-cond-quad} with the Gaussian prior gives a quadratic log density
whose precision and precision-weighted mean are
$\rho^{-2}+\sum_t\mathbb E_q[\omega_t]$ and
$\rho^{-2}\nu+\sum_t\kappa_t$, respectively. This proves the stated Gaussian update. Similarly,
$\log q^*(\omega_t)=\mathbb E_{q(\theta)}[\log p(y_t,\omega_t\mid\theta)]+\mathrm{const}$
is the log density of an exponentially tilted $\mathrm{PG}(b_t,0)$ variable with tilt
$c_t=\sqrt{\mathbb E_q[\theta^2]}$, so $q^*(\omega_t)=\mathrm{PG}(b_t,c_t)$. Its standard mean
is $b_t\tanh(c_t/2)/(2c_t)$, with the continuous limit $b_t/4$ at $c_t=0$. Each update is the
exact maximizer of the same block ELBO with the other factors fixed; therefore the ELBO is
nondecreasing after every coordinate update.
\end{proof}

\begin{algorithm}[H]
\caption{GLMBlock\_PG\_VB$(y,m,i,j,\nu,\rho^2)$ \quad (Gaussian prior, PG mean-field)}
\label{alg:pg-vb}
\KwIn{Binomial/logistic block with trial counts $m_t$; Gaussian prior $\mathcal{N}(\nu,\rho^2)$}
Initialize $q(\omega_t)$ and $q(\theta)=\mathcal{N}(\mu,\sigma^2)$\;
\Repeat{ELBO converged}{
  Set $b_t\leftarrow w_t m_t$ and $\kappa_t\leftarrow w_t(y_t-m_t/2)$ if likelihood-power weights are used; otherwise $w_t=1$\;
  $\sigma^{-2}\leftarrow \rho^{-2} + \sum_{t=i+1}^j \mathbb{E}_q[\omega_t]$\;
  $\mu \leftarrow \sigma^2 \left(\rho^{-2}\nu + \sum_{t=i+1}^j \kappa_t\right)$\;
  \ForEach{$t$}{
    Update $q(\omega_t)=\mathrm{PG}(w_t m_t,c_t)$ using current $\mathbb{E}_q[\theta^2]$\;
  }
}
Return a lower bound $\underline{A}^{(0)}_{ij}$ and approximate moments from $q(\theta)$.
\end{algorithm}

\paragraph{Using approximate block evidences in BayesBreak.}
Once any block-evidence approximation is available (Laplace, JJ, EP, PG-VB, or another surrogate), BayesBreak uses it exactly as an input score array: build a triangular array of approximate block log-evidences and run the same DP recursions. The resulting posterior summaries are exact for that approximate score array, not automatically exact for the original non-conjugate model.
Algorithm~\ref{alg:precompute-nonconj} provides a generic recipe for constructing this array in a model-agnostic way, emphasizing reuse of block-local optimizers and cached sufficient-statistic summaries.

\begin{algorithm}[H]
\caption{Precompute block evidences for non-conjugate GLMs}
\label{alg:precompute-nonconj}
\KwIn{$\{(x_t,y_t,w_t)\}_{t=1}^n$, method $\in\{$Laplace, JJ, PG--VB, EP$\}$}
\For{$i=0$ \KwTo $n-1$}{
  Initialize running summaries (as needed)\;
  \For{$j=i+1$ \KwTo $n$}{
    Update block summaries for $(i,j]$\;
    Call chosen block routine to compute $\widehat{A}^{(0)}_{ij}$ (and optional moments)\;
  }
}
Run DP (Algorithm~\ref{alg:dp-core}) with resulting block evidences.
\end{algorithm}

\subsection{Propagation of block approximation error through the DP}

\begin{assumption}[Uniform per-block log-evidence error]
\label{ass:uniform-block-error}
Let $\widehat A^{(0)}_{ij}$ denote the approximate block evidence produced by the chosen
non-conjugate block routine (one of Laplace, JJ, EP, PG mean field, or 1-D quadrature) on
candidate block $(i,j]$, and set $\delta_{ij}:=\log\widehat A^{(0)}_{ij}-\log A^{(0)}_{ij}$.
There exists $\varepsilon>0$ such that, for every candidate block $(i,j]$ with
$0\le i<j\le n$ reachable by some $k$-segmentation with $k\le k_{\max}$,
\[
|\delta_{ij}|\le \varepsilon.
\]
\end{assumption}

Assumption~\ref{ass:uniform-block-error} is an additional numerical-analysis requirement, not a
consequence of naming an approximation routine. It must be established independently for the
chosen family, parameter domain, candidate-block support, transformation, initialization,
convergence event, and tail treatment. Nested quadrature, a high-accuracy reference calculation,
or a method-specific remainder theorem can provide such a certificate in a particular analysis.
Convergence of an optimizer, monotonicity of an ELBO, or an EP moment match alone does not provide
the required two-sided uniform log-evidence bound.

\begin{proposition}[Log-evidence and posterior-odds stability]
\label{prop:stability}
Under Assumption~\ref{ass:uniform-block-error}, for every DP state $(k,j)$ with $k\le k_{\max}$,
\[
e^{-k\varepsilon}L_{k j}\le \widehat{L}_{k j}\le e^{k\varepsilon}L_{k j},
\qquad
e^{-k\varepsilon}R_{k j}\le \widehat{R}_{k j}\le e^{k\varepsilon}R_{k j}.
\]
Consequently:
\begin{enumerate}
\item For any $k,k'\le k_{\max}$, the posterior odds for the segment count satisfy
\begin{equation}
\left|
\log\frac{\widehat{P}(k\mid y)}{\widehat{P}(k'\mid y)}
-\log\frac{P(k\mid y)}{P(k'\mid y)}
\right|
\le (k+k')\varepsilon.
\label{eq:stability-k-odds}
\end{equation}
\item For any fixed $k\le k_{\max}$ and any two candidate locations $h,h'$ of boundary $t_p$,
\begin{equation}
\left|
\log\frac{\widehat{P}(t_p=h\mid y,k)}{\widehat{P}(t_p=h'\mid y,k)}
-\log\frac{P(t_p=h\mid y,k)}{P(t_p=h'\mid y,k)}
\right|
\le 2k\varepsilon.
\label{eq:stability-b-odds}
\end{equation}
\end{enumerate}
Thus uniform blockwise log-evidence control yields uniform control of global posterior \emph{odds}. Coarse absolute probability bounds follow by normalizing finite weights, as in Corollary~\ref{cor:probability-error-conversion}; a margin condition is needed instead for preservation of a modal ranking.
\end{proposition}

\begin{proof}
Any term contributing to $L_{k j}$ is a product of exactly $k$ block evidences. Hence the corresponding
approximate term differs by a multiplicative factor in $[e^{-k\varepsilon},e^{k\varepsilon}]$. Summing over all
paths yields the same sandwich bound for $\widehat{L}_{k j}$, and the suffix bound is identical. Equation
\eqref{eq:stability-k-odds} follows immediately because $P(k\mid y)$ is proportional to $p(k)L_{k n}/C_k$ and the
prior and normalization constants are unchanged. For boundary odds, note that
$P(t_p=h\mid y,k)\propto L_{p,h}R_{k-p,h}$. The numerator for each candidate location involves $k$ block factors in
total, so each candidate-specific weight is perturbed by at most $e^{\pm k\varepsilon}$; comparing two candidates
therefore yields the factor $e^{\pm 2k\varepsilon}$ and hence \eqref{eq:stability-b-odds}.
\end{proof}

\begin{corollary}[Probability-error conversion for normalized finite weights]
\label{cor:probability-error-conversion}
Let $p_a=w_a/\sum_b w_b$ and $\widehat p_a=\widehat w_a/\sum_b\widehat w_b$ be two
finite normalized distributions with $|\log\widehat w_a-\log w_a|\le\eta$ for every state $a$. Then
$e^{-2\eta}\le \widehat p_a/p_a\le e^{2\eta}$ for every $a$, and
$\|\widehat p-p\|_{\mathrm{TV}}\le \min\{1,e^{2\eta}-1\}$. Consequently, Assumption~\ref{ass:uniform-block-error}
implies a fixed-$k$ boundary-distribution total-variation bound with $\eta=k\varepsilon$, and a
segment-count posterior bound with $\eta=k_{\max}\varepsilon$. These absolute probability bounds are
coarser than the odds bounds but require no additional margin condition.
\end{corollary}

\begin{proof}
The normalizing denominator $\sum_b p_b\exp(\delta_b)$ lies in $[e^{-\eta},e^{\eta}]$ when
$\delta_a=\log\widehat w_a-\log w_a$. Hence each probability ratio lies in
$[e^{-2\eta},e^{2\eta}]$. Summing the positive parts of $\widehat p_a-p_a$ gives the exponential bound, and total variation is always at most one.
\end{proof}

\begin{corollary}[Preservation of boundary rankings under a margin condition]
\label{cor:ranking}
Under the hypotheses of Proposition~\ref{prop:stability}, if there is a \emph{log-odds margin}
$\Delta>0$ separating the best candidate $h^\star$ from every alternative,
\[
\log\frac{P(t_p=h^\star\mid y,k)}{P(t_p=h\mid y,k)}\ge \Delta
\quad\text{for all }h\ne h^\star,
\]
then for every $\varepsilon<\Delta/(2k)$ the approximate posterior preserves the exact modal
boundary location: $\arg\max_h \widehat P(t_p=h\mid y,k)=h^\star$.
\end{corollary}

\begin{proof}
By \eqref{eq:stability-b-odds}, the log-odds between $h^\star$ and any alternative $h$ is
perturbed by at most $2k\varepsilon<\Delta$, so the exact log-odds positivity is preserved.
\end{proof}


\begin{proofobligation}[Routine-specific block-error certification]
\label{prop:uniform-bounds}
No routine-wide uniform rate is asserted for Laplace, Jaakkola--Jordan, expectation propagation,
P\'olya--Gamma mean field, or adaptive quadrature without assumptions tailored to the family,
parameter domain, block support, initialization, convergence event, and tails. A valid run must
instead provide a reachable-block certificate that separates optimization residual, truncation or
tail error, quadrature error, and comparison-to-reference error. The conditional global theorem
below consumes the maximum certified log-score error; it does not manufacture that certificate.
\end{proofobligation}



\begin{remark}[A small numerical illustration of the stability bound]
\label{rem:stability-numerical}
Consider $n=500$, $k_{\max}=10$, and a non-conjugate block routine whose verified uniform
per-block log-evidence error over all reachable blocks is $\varepsilon=0.02$. Proposition~\ref{prop:stability}
then guarantees segment-count log-odds error at most $(5+10)\times0.02=0.30$ when comparing
$k=5$ to $k=10$, and boundary log-odds error at most $2\cdot10\cdot0.02=0.40$ under any fixed
$k\le10$. In terms of probability ratios, these correspond to factors of about $e^{0.30}$ and
$e^{0.40}$. The numbers in this illustration are conditional on having already verified the
uniform block-error bound; they are not a generic guarantee for any chosen approximation routine.
\end{remark}

\begin{remark}[When each non-conjugate routine struggles]
\label{rem:failure-modes}
The stability bound in Proposition~\ref{prop:stability} is tight only when the uniform blockwise
error $\varepsilon$ is actually small. Each non-conjugate routine has its own failure mode:
Laplace (Theorem~\ref{thm:laplace}) requires posterior concentration near $\hat\theta_{ij}$, so
it can be inaccurate on very short segments ($W_{ij}$ small) or for heavy-tailed priors;
mean-field variational Bayes (Proposition~\ref{prop:pg-mf}) can be overconfident and
underestimate posterior variance, biasing moment numerators; EP can fail to converge for
non-log-concave likelihoods; 1D quadrature is accurate for smooth 1D integrands but expensive
in higher dimensions. For these reasons, monitor $\varepsilon$ empirically per
dataset rather than assuming it. A practical validation checklist is to compare a subset of blocks against
quadrature or another trusted reference, report both maximum and high-quantile block errors, inspect
errors along the exported MAP path, and assess sensitivity of $P(k\mid y)$ and boundary marginals.
\end{remark}

\paragraph{Approximation-validation checklist.}
Before interpreting non-conjugate posterior summaries, compare a subset of block evidences against
higher-accuracy quadrature or another trusted reference; report maximum and high-quantile block
errors over reachable blocks; inspect errors on the MAP path; and assess sensitivity of $p(k\mid y)$,
boundary marginals, and exported segmentations to the approximation method. The ranking corollary
is meaningful only when the exact posterior margin is larger than the propagated approximation error.


\begin{figure}[H]
\centering
\resizebox{0.96\textwidth}{!}{\begin{tikzpicture}[node distance=7mm and 8mm]
\node[secondarybox,text width=33mm] (ref) {Reference and approximate log segment marginal likelihoods};
\node[primarybox,right=of ref,text width=29mm] (eps) {$|\widehat\ell_{ij}-\ell_{ij}|\le\varepsilon$ for every admissible segment};
\node[primarybox,right=of eps,text width=31mm] (part) {$k$-segment partition log-weight error at most $k\varepsilon$};
\node[successbox,right=of part,text width=33mm] (global) {Bounds for marginal likelihoods, posterior odds, and modal changepoints};
\draw[arrPrimary] (ref)--(eps);
\draw[arrPrimary] (eps)--(part);
\draw[arrPrimary] (part)--(global);
\node[below=5mm of part,font=\scriptsize,text width=66mm,align=center] {The implication is conditional: the numerical method must establish the segmentwise bound under its own assumptions.};
\end{tikzpicture}
}
\caption{Conditional propagation from a reachable-block log-score tolerance to global evidence and posterior-odds bounds. The local numerical routine remains responsible for establishing the tolerance.}
\label{fig:local-global-stability}
\end{figure}
